\section{Ethics and Responsible Disclosure}
% provenance: mis_01KXT2D8VVC1CDNT9WJ4CFGH9P scope boundary (synthetic keys only; no production testing; no vendor-bound exploit tooling).

This work studies the consequences of credential compromise in order to bound them, and it
carries a corresponding responsibility. All experiments use synthetic and lab-generated
certificates. We use no real vendor keys or secrets, we test against no production DER fleet
or any internet-reachable device, and we build no exploit tooling bound to a specific vendor
product. The malicious dispatch programs exist only inside the feeder co-simulation and are
never emitted to hardware.

The physical-consequence results describe what a compromised but protocol-conformant
credential can already do under the deployed standard; they do not introduce a new
vulnerability so much as measure an existing exposure that the standard's prohibition on
revocation leaves unmitigated. Where results identify a concrete deployment weakness, we will
follow coordinated disclosure with the relevant standards body and affected implementers
before public release, and we will time any artifact release accordingly. The redesign and
its migration path are offered as the constructive counterpart to the measurement: the point
of quantifying the blast radius is to shrink it.
